\documentclass[12pt, titlepage]{article}

\usepackage{amsmath, mathtools}

\usepackage[numbers,square,sort&compress]{natbib}
\usepackage{amsfonts}
\usepackage{amssymb}
\usepackage{graphicx}
\usepackage{colortbl}
\usepackage{xr}
\usepackage{hyperref}
\usepackage{longtable}
\usepackage{xfrac}
\usepackage{tabularx}
\usepackage{float}
\usepackage{siunitx}
\usepackage{booktabs}
\usepackage{multirow}
\usepackage[section]{placeins}
\usepackage{caption}
\usepackage{fullpage}

\hypersetup{
bookmarks=true,     % show bookmarks bar?
colorlinks=true,       % false: boxed links; true: colored links
linkcolor=red,          % color of internal links (change box color with linkbordercolor)
citecolor=blue,      % color of links to bibliography
filecolor=magenta,  % color of file links
urlcolor=cyan          % color of external links
}

\usepackage{array}

\externaldocument[SRS-]{../../SRS-IEEE/SRS}
\externaldocument[MG-]{../SoftArchitecture/MG}

\newcommand{\mref}[1]{M\ref{MG-#1}}
\newcommand{\frref}[1]{FR\ref{SRS-#1}}
\newcommand{\nfrref}[1]{NFR\ref{SRS-#1}}
\newcommand{\drref}[1]{DR\ref{SRS-#1}}

\input{../../Comments.text}
\input{../../Common.text}

% IEEEtran.bst control: write repeated author names in full instead of a dash.
% \bstctlcite is defined by the IEEEtran class; article needs it here.
\makeatletter
\def\bstctlcite{\@ifnextchar[{\@bstctlcite}{\@bstctlcite[@auxout]}}
\def\@bstctlcite[#1]#2{\@bsphack\@for\@citeb:=#2\do{%
  \edef\@citeb{\expandafter\@firstofone\@citeb}%
  \if@filesw\immediate\write\csname #1\endcsname{\string\citation{\@citeb}}\fi}%
  \@esphack}
\makeatother

\begin{document}
\bstctlcite{IEEEexample:BSTcontrol}

\title{Module Interface Specification for \progname{}}

\author{\authname}

\date{\today}

\maketitle

\pagenumbering{roman}

\section{Revision History}

\begin{tabularx}{\textwidth}{p{3cm}p{2cm}X}
\toprule {\bf Date} & {\bf Version} & {\bf Notes}\\
\midrule
Nov.\ 16, 2026 & 1.0 & Initial Version\\
\bottomrule
\end{tabularx}

~\newpage

\section{Symbols, Abbreviations and Acronyms}

Symbols, abbreviations and acronyms defined in the SRS~\citep{TeamBVSH2026SRS}
(for example AIR, CAS, CEAB, DR, FR, GA, LLM, NFR and PDF) and in the
MG~\citep{TeamBVSH2026MG} (for example AC, M and UC) keep their meaning there.
This document adds the following:

\renewcommand{\arraystretch}{1.2}
\begin{tabular}{l l}
  \toprule
  \textbf{symbol} & \textbf{description}\\
  \midrule
  MIS & Module Interface Specification\\
  $\mathbb{B}$ & Boolean: true or false\\
  -- & No input, no output or no exception\\
  \bottomrule
\end{tabular}

\newpage

\tableofcontents

\newpage

\pagenumbering{arabic}

\section{Introduction}

%The following document details the Module Interface Specifications for
%\wss{Fill in your project name and description}

%Complementary documents include the System Requirement Specifications
%and Module Guide.  The full documentation and implementation can be
%found at \url{...}.  \wss{provide the url for your repo}

The following document details the Module Interface Specifications for the ASSURE Project. 
The system collects assessment information from CAS course outlines, shows both confirmed 
and tentative assessment dates in a shared calendar, shows how many assessments each 
student cohort has per week, and checks each outline against McMaster's course 
outline policy~\citep{McMasterUCMP2025} and the CEAB graduate attributes~\citep{CEAB2025}.\\

\noindent Complementary documents include the System Requirement Specifications and Module Guide.
The full documentation and implementation can be found at
\url{https://github.com/shivyam/cas-course-optimization}.


\section{Notation}

The structure of the MIS for modules comes from Hoffman and Strooper~\citep{HoffmanAndStrooper1995},
with the addition that template modules have been adapted from
Ghezzi et al.~\citep{GhezziEtAl2003}.  The mathematical notation comes from Chapter 3 of
Hoffman and Strooper~\citep{HoffmanAndStrooper1995}.  For instance, the symbol := is used for a
multiple assignment statement and conditional rules follow the form $(c_1
\Rightarrow r_1 | c_2 \Rightarrow r_2 | ... | c_n \Rightarrow r_n )$.

The following table summarizes the primitive data types used by \progname. 

\begin{center}
\renewcommand{\arraystretch}{1.2}
\noindent 
\begin{tabular}{l l p{7.5cm}} 
\toprule 
\textbf{Data Type} & \textbf{Notation} & \textbf{Description}\\ 
\midrule
boolean & $\mathbb{B}$ & true or false\\
character & char & a single symbol or digit\\
integer & $\mathbb{Z}$ & a number without a fractional component in (-$\infty$, $\infty$) \\
natural number & $\mathbb{N}$ & a number without a fractional component in [1, $\infty$) \\
real & $\mathbb{R}$ & any number in (-$\infty$, $\infty$)\\
\bottomrule
\end{tabular} 
\end{center}

\noindent
The specification of \progname \ uses some derived data types: sequences, strings, and
tuples. Sequences are lists filled with elements of the same data type. Strings
are sequences of characters. Tuples contain a list of values, potentially of
different types. In addition, \progname \ uses functions, which
are defined by the data types of their inputs and outputs. Local functions are
described by giving their type signature followed by their specification.

The following abstract types are used by the access programs. Their fields
follow the data dictionary of the SRS~\citep{TeamBVSH2026SRS}; how they are
stored is the secret of the Data Store (\mref{mStore}).

\begin{center}
\renewcommand{\arraystretch}{1.2}
\begin{longtable}{>{\raggedright\arraybackslash}p{4.6cm} >{\raggedright\arraybackslash}p{10.1cm}}
\toprule
\textbf{Type} & \textbf{Description}\\
\midrule
CourseID, OutlineID, AssessmentID, FindingID, RecID, ValueID, ID & Identifiers given out by the Data Store (\mref{mStore})\\
User & A user and their role in \{administrator, instructor\}\\
Session & A logged-in session\\
Credentials & What a user gives to log in\\
File & An uploaded or generated file and its name\\
Outline & A submitted course outline and the course it belongs to\\
Value & One extracted value, with its source passage and page\\
Assessment & (name, type, date or missing, weight, tentative: $\mathbb{B}$)\\
CourseRecord & The assessments and learning outcomes extracted from one outline\\
Finding & (course, kind, message, source passage, accepted: $\mathbb{B}$)\\
AuditEntry & (user, time, value or finding changed, previous value, reason)\\
DateRecommendation & (assessment, suggested date, explanation, decision)\\
CourseCode & A course code, such as SFWRENG 4X03\\
Cohort & (program, year level), for example (COMPSCI, 2)\\
Date, Week, Term & A calendar date; a week of term; a term, such as Winter 2027\\
Status & (missing values: $\mathbb{N}$, open findings: $\mathbb{N}$, unconfirmed values: $\mathbb{N}$)\\
CalendarFilter & The courses, year levels and programs to show\\
Request, Response & A request from the Web UI (\mref{mUI}) and the reply to it\\
Entity, EntityKind, Filter & A stored record, its kind, and a condition on its fields\\
\bottomrule
\end{longtable}
\end{center}

\section{Module Decomposition}

The following table is taken directly from the Module
Guide~\citep{TeamBVSH2026MG}. The Hardware-Hiding module (\mref{mHH}) is
provided by the operating system and has no MIS.

\begin{table}[h!]
\centering
\begin{tabular}{p{0.3\textwidth} p{0.6\textwidth}}
\toprule
\textbf{Level 1} & \textbf{Level 2}\\
\midrule

{Hardware-Hiding} & Hardware-Hiding (\mref{mHH})\\
\midrule

\multirow{16}{0.3\textwidth}{Behaviour-Hiding}
& Outline Intake (\mref{mIntake})\\
& Extraction (\mref{mExtract})\\
& Extraction Review (\mref{mReview})\\
& Assessment Schedule (\mref{mSchedule})\\
& Workload Analysis (\mref{mWorkload})\\
& Policy Compliance (\mref{mPolicy})\\
& Graduate Attribute Check (\mref{mGA})\\
& Findings and Override (\mref{mFindings})\\
& Dashboard (\mref{mDashboard})\\
& Date Suggestion (\mref{mSuggest})\\
& Export (\mref{mExport})\\
& Access Control (\mref{mAccess})\\
& Web UI (\mref{mUI})\\
& Request Controller (\mref{mControl})\\
& Cohort (\mref{mCohort})\\
& Academic Calendar (\mref{mCalendar})\\
\midrule

\multirow{3}{0.3\textwidth}{Software Decision}
& Document Reader (\mref{mReader})\\
& LLM Client (\mref{mLLM})\\
& Data Store (\mref{mStore})\\
\bottomrule

\end{tabular}
\caption{Module Hierarchy}
\label{TblMH}
\end{table}

\newpage
~\newpage

\section{MIS of Outline Intake Module (\mref{mIntake})} \label{mis-mIntake}

\subsection{Module}

OutlineIntake

\subsection{Uses}

\mref{mExtract}, \mref{mReader}, \mref{mStore}

\subsection{Syntax}

\subsubsection{Exported Constants}

\begin{itemize}
\item \textit{MAX\_UPLOAD\_SIZE}: Largest outline file accepted, 5 MB.
\end{itemize}

\subsubsection{Exported Access Programs}

\begin{center}\small
\begin{tabular}{>{\raggedright\arraybackslash}p{3.3cm} >{\raggedright\arraybackslash}p{4.1cm} >{\raggedright\arraybackslash}p{3.4cm} >{\raggedright\arraybackslash}p{4.0cm}}
\hline
\textbf{Name} & \textbf{In} & \textbf{Out} & \textbf{Exceptions} \\
\hline
submit & files: seq of File, u: User & seq of OutlineID & UnsupportedFormat, FileTooLarge, EmptyFile, NotAuthorized \\
getOutline & id: OutlineID & Outline & NotFound \\
\hline
\end{tabular}
\end{center}

\subsection{Semantics}

\subsubsection{State Variables}

% None.
\todo{Rev 0}

\subsubsection{Environment Variables}

% None.
\todo{Rev 0}


\subsubsection{Assumptions}

% Each submitted file is one course outline.
\todo{Rev 0}


\subsubsection{Access Routine Semantics}

% \todo{Rev 0: transition, output and exception for each access program.}
\todo{Rev 0}

\subsubsection{Local Functions}

\todo{Rev 0}

\newpage

\section{MIS of Extraction Module (\mref{mExtract})} \label{mis-mExtract}

\subsection{Module}

Extraction

\subsection{Uses}

\mref{mReader}, \mref{mLLM}, \mref{mStore}

\subsection{Syntax}

\subsubsection{Exported Constants}

None.


\subsubsection{Exported Access Programs}

\begin{center}\small
\begin{tabular}{>{\raggedright\arraybackslash}p{3.3cm} >{\raggedright\arraybackslash}p{4.1cm} >{\raggedright\arraybackslash}p{3.4cm} >{\raggedright\arraybackslash}p{4.0cm}}
\hline
\textbf{Name} & \textbf{In} & \textbf{Out} & \textbf{Exceptions} \\
\hline
extract & id: OutlineID & CourseRecord & ModelUnavailable, ExtractionTimeout, InvalidModelOutput \\
missingFields & c: CourseID & seq of Finding & NotFound \\
\hline
\end{tabular}
\end{center}

\subsection{Semantics}

\subsubsection{State Variables}

% None.
\todo{Rev 0}


\subsubsection{Environment Variables}

% None.
\todo{Rev 0}


\subsubsection{Assumptions}

% The text returned by the Document Reader (\mref{mReader}) is in reading order.
\todo{Rev 0}

\subsubsection{Access Routine Semantics}

% \todo{Rev 0: transition, output and exception for each access program.}
\todo{Rev 0}


\subsubsection{Local Functions}

\todo{Rev 0}


\newpage

\section{MIS of Extraction Review Module (\mref{mReview})} \label{mis-mReview}

\subsection{Module}

ExtractionReview

\subsection{Uses}

\mref{mStore}

\subsection{Syntax}

\subsubsection{Exported Constants}

None.

\subsubsection{Exported Access Programs}

\begin{center}\small
\begin{tabular}{>{\raggedright\arraybackslash}p{3.3cm} >{\raggedright\arraybackslash}p{4.1cm} >{\raggedright\arraybackslash}p{3.4cm} >{\raggedright\arraybackslash}p{4.0cm}}
\hline
\textbf{Name} & \textbf{In} & \textbf{Out} & \textbf{Exceptions} \\
\hline
getForReview & c: CourseID & CourseRecord & NotFound \\
correct & v: ValueID, x: Value, u: User & -- & NotFound, NotAuthorized \\
confirm & c: CourseID, u: User & -- & NotAuthorized, UnreviewedValues \\
isConfirmed & c: CourseID & $\mathbb{B}$ & NotFound \\
\hline
\end{tabular}
\end{center}

\subsection{Semantics}

\subsubsection{State Variables}

% \begin{itemize}
% \item reviewStatus: ValueID $\rightarrow$ \{extracted, corrected, confirmed\}
% \end{itemize}
% Stored through the Data Store (\mref{mStore}).
\todo{Rev 0}


\subsubsection{Environment Variables}

% None.
\todo{Rev 0}

\subsubsection{Assumptions}

% None.
\todo{Rev 0}


\subsubsection{Access Routine Semantics}

% \todo{Rev 0: transition, output and exception for each access program.}
\todo{Rev 0}


\subsubsection{Local Functions}

\todo{Rev 0}

\newpage

\section{MIS of Assessment Schedule Module (\mref{mSchedule})} \label{mis-mSchedule}

\subsection{Module}

AssessmentSchedule

\subsection{Uses}

\mref{mReview}, \mref{mWorkload}, \mref{mPolicy}, \mref{mStore}

\subsection{Syntax}

\subsubsection{Exported Constants}

None.

\subsubsection{Exported Access Programs}

\begin{center}\small
\begin{tabular}{>{\raggedright\arraybackslash}p{3.3cm} >{\raggedright\arraybackslash}p{4.1cm} >{\raggedright\arraybackslash}p{3.4cm} >{\raggedright\arraybackslash}p{4.0cm}}
\hline
\textbf{Name} & \textbf{In} & \textbf{Out} & \textbf{Exceptions} \\
\hline
setTentativeDate & a: AssessmentID, d: Date, u: User & -- & NotOwner, DateOutsideTerm \\
approveDate & a: AssessmentID, d: Date, u: User & -- & NotOwner \\
getSchedule & c: CourseID & seq of Assessment & NotFound \\
\hline
\end{tabular}
\end{center}

\subsection{Semantics}

\subsubsection{State Variables}

% \begin{itemize}
% \item tentative: AssessmentID $\rightarrow$ Date
% \end{itemize}
% Stored through the Data Store (\mref{mStore}).
\todo{Rev 0}

\subsubsection{Environment Variables}

% None.
\todo{Rev 0}


\subsubsection{Assumptions}

% Only the instructor of a course changes its dates (\uref{ucAuthority}).
\todo{Rev 0}

\subsubsection{Access Routine Semantics}

% \todo{Rev 0: transition, output and exception for each access program.}
\todo{Rev 0}


\subsubsection{Local Functions}

% \todo{Rev 0, as needed.}
\todo{Rev 0}

\newpage

% Section 10
\section{MIS of Workload Analysis Module (\mref{mWorkload})} \label{mis-mWorkload}

\subsection{Module}
WorkloadAnalysis

\subsection{Uses}
\mref{mReview}, \mref{mStore}, \mref{mCohort}, \mref{mCalendar}

\subsection{Syntax}

\subsubsection{Exported Constants}
None.

\subsubsection{Exported Access Programs}
\begin{center}\small
\begin{tabular}{>{\raggedright\arraybackslash}p{3.3cm} >{\raggedright\arraybackslash}p{4.1cm} >{\raggedright\arraybackslash}p{3.4cm} >{\raggedright\arraybackslash}p{4.0cm}}
\hline
\textbf{Name} & \textbf{In} & \textbf{Out} & \textbf{Exceptions} \\
\hline
weeklyDensity & k: Cohort, t: Term & seq of (Week, $\mathbb{N}$) & UnknownCohort \\
restrictedConflicts & c: CourseID & seq of Finding & NotFound \\
reanalyze & c: CourseID & -- & NotFound \\
\hline
\end{tabular}
\end{center}

\subsection{Semantics}

\subsubsection{State Variables}
\todo{Rev 0}

\subsubsection{Environment Variables}
\todo{Rev 0}

\subsubsection{Assumptions}
\todo{Rev 0}

\subsubsection{Access Routine Semantics}
\todo{Rev 0}

\subsubsection{Local Functions}
\todo{Rev 0}

\newpage

% Section 11
\section{MIS of Policy Compliance Module (\mref{mPolicy})} \label{mis-mPolicy}

\subsection{Module}
PolicyCompliance

\subsection{Uses}
\mref{mReview}, \mref{mFindings}, \mref{mStore}, \mref{mCalendar}

\subsection{Syntax}

\subsubsection{Exported Constants}
\begin{itemize}
\item \textit{WEIGHT\_TOTAL}: policy value from SRS Appendix C.
\item \textit{MAX\_SINGLE\_WEIGHT}: policy value from SRS Appendix C.
\item \textit{BAN\_CLASS\_DAYS}: policy value from SRS Appendix C.
\item \textit{BAN\_WEIGHT\_LIMIT}: policy value from SRS Appendix C.
\item \textit{EARLY\_FEEDBACK\_L2}: policy value from SRS Appendix C.
\item \textit{EARLY\_FEEDBACK\_L3}: policy value from SRS Appendix C.
\item \textit{OUTLINE\_LEAD\_TIME}: policy value from SRS Appendix C.
\end{itemize}

\subsubsection{Exported Access Programs}
\begin{center}\small
\begin{tabular}{>{\raggedright\arraybackslash}p{3.3cm} >{\raggedright\arraybackslash}p{4.1cm} >{\raggedright\arraybackslash}p{3.4cm} >{\raggedright\arraybackslash}p{4.0cm}}
\hline
\textbf{Name} & \textbf{In} & \textbf{Out} & \textbf{Exceptions} \\
\hline
checkPolicy & c: CourseID & seq of Finding & NotConfirmed, NotFound \\
\hline
\end{tabular}
\end{center}

\subsection{Semantics}

\subsubsection{State Variables}
\todo{Rev 0}

\subsubsection{Environment Variables}
\todo{Rev 0}

\subsubsection{Assumptions}
\todo{Rev 0}

\subsubsection{Access Routine Semantics}
\todo{Rev 0}

\subsubsection{Local Functions}
\todo{Rev 0}

\newpage

% Section 12
\section{MIS of Graduate Attribute Check Module (\mref{mGA})} \label{mis-mGA}

\subsection{Module}
GACheck

\subsection{Uses}
\mref{mReview}, \mref{mFindings}, \mref{mStore}

\subsection{Syntax}

\subsubsection{Exported Constants}
\begin{itemize}
\item \textit{GA\_LIST}: The twelve CEAB graduate attributes~\citep{CEAB2025}.
\end{itemize}

\subsubsection{Exported Access Programs}
\begin{center}\small
\begin{tabular}{>{\raggedright\arraybackslash}p{3.3cm} >{\raggedright\arraybackslash}p{4.1cm} >{\raggedright\arraybackslash}p{3.4cm} >{\raggedright\arraybackslash}p{4.0cm}}
\hline
\textbf{Name} & \textbf{In} & \textbf{Out} & \textbf{Exceptions} \\
\hline
checkGA & c: CourseID & seq of Finding & NotConfirmed, NoMasterRecord \\
\hline
\end{tabular}
\end{center}

\subsection{Semantics}

\subsubsection{State Variables}
\todo{Rev 0}

\subsubsection{Environment Variables}
\todo{Rev 0}

\subsubsection{Assumptions}
\todo{Rev 0}

\subsubsection{Access Routine Semantics}
\todo{Rev 0}

\subsubsection{Local Functions}
\todo{Rev 0}

\newpage

% Section 13
\section{MIS of Findings and Override Module (\mref{mFindings})} \label{mis-mFindings}

\subsection{Module}
Findings

\subsection{Uses}
\mref{mStore}

\subsection{Syntax}

\subsubsection{Exported Constants}
None.

\subsubsection{Exported Access Programs}
\begin{center}\small
\begin{tabular}{>{\raggedright\arraybackslash}p{3.3cm} >{\raggedright\arraybackslash}p{4.1cm} >{\raggedright\arraybackslash}p{3.4cm} >{\raggedright\arraybackslash}p{4.0cm}}
\hline
\textbf{Name} & \textbf{In} & \textbf{Out} & \textbf{Exceptions} \\
\hline
record & f: Finding & FindingID & -- \\
findings & c: CourseID & seq of Finding & NotFound \\
override & f: FindingID, r: String, u: User & -- & NotAdmin, EmptyReason, NotFound \\
auditTrail & c: CourseID & seq of AuditEntry & NotAdmin \\
\hline
\end{tabular}
\end{center}

\subsection{Semantics}

\subsubsection{State Variables}
\todo{Rev 0}

\subsubsection{Environment Variables}
\todo{Rev 0}

\subsubsection{Assumptions}
\todo{Rev 0}

\subsubsection{Access Routine Semantics}
\todo{Rev 0}

\subsubsection{Local Functions}
\todo{Rev 0}

\newpage

% Section 14
\section{MIS of Dashboard Module (\mref{mDashboard})} \label{mis-mDashboard}

\subsection{Module}
Dashboard

\subsection{Uses}
\mref{mReview}, \mref{mSchedule}, \mref{mFindings}, \mref{mStore}

\subsection{Syntax}

\subsubsection{Exported Constants}
None.

\subsubsection{Exported Access Programs}
\begin{center}\small
\begin{tabular}{>{\raggedright\arraybackslash}p{3.3cm} >{\raggedright\arraybackslash}p{4.1cm} >{\raggedright\arraybackslash}p{3.4cm} >{\raggedright\arraybackslash}p{4.0cm}}
\hline
\textbf{Name} & \textbf{In} & \textbf{Out} & \textbf{Exceptions} \\
\hline
visibleCourses & u: User & seq of CourseID & -- \\
courseStatus & c: CourseID & Status & NotFound \\
termCompleteness & t: Term & $\mathbb{R}$ (share of courses fully planned) & -- \\
\hline
\end{tabular}
\end{center}

\subsection{Semantics}

\subsubsection{State Variables}
\todo{Rev 0}

\subsubsection{Environment Variables}
\todo{Rev 0}

\subsubsection{Assumptions}
\todo{Rev 0}

\subsubsection{Access Routine Semantics}
\todo{Rev 0}

\subsubsection{Local Functions}
\todo{Rev 0}

\newpage

% Section 15
\section{MIS of Date Suggestion Module (\mref{mSuggest})} \label{mis-mSuggest}

\subsection{Module}
DateSuggestion

\subsection{Uses}
\mref{mWorkload}, \mref{mStore}, \mref{mCalendar}

\subsection{Syntax}

\subsubsection{Exported Constants}
\begin{itemize}
\item \textit{MIN\_ALTERNATIVES}: Fewest dates offered per request, from SRS Appendix C.
\end{itemize}

\subsubsection{Exported Access Programs}
\begin{center}\small
\begin{tabular}{>{\raggedright\arraybackslash}p{3.3cm} >{\raggedright\arraybackslash}p{4.1cm} >{\raggedright\arraybackslash}p{3.4cm} >{\raggedright\arraybackslash}p{4.0cm}}
\hline
\textbf{Name} & \textbf{In} & \textbf{Out} & \textbf{Exceptions} \\
\hline
suggest & a: AssessmentID, u: User & seq of DateRecommendation & NotOwner, NoConflict \\
decide & r: RecID, accept: $\mathbb{B}$, u: User & -- & NotOwner, StaleSuggestion \\
\hline
\end{tabular}
\end{center}

\subsection{Semantics}

\subsubsection{State Variables}
\todo{Rev 0}

\subsubsection{Environment Variables}
\todo{Rev 0}

\subsubsection{Assumptions}
\todo{Rev 0}

\subsubsection{Access Routine Semantics}
\todo{Rev 0}

\subsubsection{Local Functions}
\todo{Rev 0}

\newpage

\section{MIS of Export Module (\mref{mExport})} \label{mis-mExport}

\subsection{Module}

Export

\subsection{Uses}

\mref{mStore}

\subsection{Syntax}

\subsubsection{Exported Constants}

None.

\subsubsection{Exported Access Programs}

\begin{center}\small
\begin{tabular}{>{\raggedright\arraybackslash}p{3.3cm} >{\raggedright\arraybackslash}p{4.1cm} >{\raggedright\arraybackslash}p{3.4cm} >{\raggedright\arraybackslash}p{4.0cm}}
\hline
\textbf{Name} & \textbf{In} & \textbf{Out} & \textbf{Exceptions} \\
\hline
exportData & t: Term, u: User & File & NotAdmin \\
\hline
\end{tabular}
\end{center}

\subsection{Semantics}

\subsubsection{State Variables}

None.

\subsubsection{Environment Variables}

None.

\subsubsection{Assumptions}

None.

\subsubsection{Access Routine Semantics}

\todo{Rev 0}

\subsubsection{Local Functions}

\todo{Rev 0}

\newpage

\section{MIS of Access Control Module (\mref{mAccess})} \label{mis-mAccess}

\subsection{Module}

AccessControl

\subsection{Uses}

\mref{mStore}

\subsection{Syntax}

\subsubsection{Exported Constants}

\begin{itemize}
\item \textit{SESSION\_TIMEOUT}: Inactivity limit before a session expires, 30 minutes (\frref{FR_Access_Timeout}).
\end{itemize}

\subsubsection{Exported Access Programs}

\begin{center}\small
\begin{tabular}{>{\raggedright\arraybackslash}p{3.3cm} >{\raggedright\arraybackslash}p{4.1cm} >{\raggedright\arraybackslash}p{3.4cm} >{\raggedright\arraybackslash}p{4.0cm}}
\hline
\textbf{Name} & \textbf{In} & \textbf{Out} & \textbf{Exceptions} \\
\hline
login & cr: Credentials & Session & InvalidCredentials \\
logout & s: Session & -- & -- \\
currentUser & s: Session & User & SessionExpired \\
canEdit & u: User, c: CourseID & $\mathbb{B}$ & -- \\
isAdmin & u: User & $\mathbb{B}$ & -- \\
\hline
\end{tabular}
\end{center}

\subsection{Semantics}

\subsubsection{State Variables}

\begin{itemize}
\item sessions: set of (Session, User, lastActive: Time)
\end{itemize}
Stored through the Data Store (\mref{mStore}).

\subsubsection{Environment Variables}

\begin{itemize}
\item clock: The server's system clock.
\end{itemize}

\subsubsection{Assumptions}

None.

\subsubsection{Access Routine Semantics}

\todo{Rev 0}

\subsubsection{Local Functions}

\todo{Rev 0}

\newpage

\section{MIS of Web UI Module (\mref{mUI})} \label{mis-mUI}

\subsection{Module}

WebUI

\subsection{Uses}

\mref{mControl}

\subsection{Syntax}

\subsubsection{Exported Constants}

None.

\subsubsection{Exported Access Programs}

\begin{center}\small
\begin{tabular}{>{\raggedright\arraybackslash}p{3.3cm} >{\raggedright\arraybackslash}p{4.1cm} >{\raggedright\arraybackslash}p{3.4cm} >{\raggedright\arraybackslash}p{4.0cm}}
\hline
\textbf{Name} & \textbf{In} & \textbf{Out} & \textbf{Exceptions} \\
\hline
showUpload & -- & -- & -- \\
showReview & c: CourseID & -- & -- \\
showCalendar & f: CalendarFilter & -- & -- \\
showFindings & c: CourseID & -- & -- \\
showDashboard & t: Term & -- & -- \\
\hline
\end{tabular}
\end{center}

\subsection{Semantics}

\subsubsection{State Variables}

None.

\subsubsection{Environment Variables}

\begin{itemize}
\item window: The user's browser window.
\end{itemize}

\subsubsection{Assumptions}

None.

\subsubsection{Access Routine Semantics}

\todo{Rev 0}

\subsubsection{Local Functions}

\todo{Rev 0}

\newpage

\section{MIS of Request Controller Module (\mref{mControl})} \label{mis-mControl}

\subsection{Module}

RequestController

\subsection{Uses}

\mref{mIntake}, \mref{mExtract}, \mref{mReview}, \mref{mSchedule}, \mref{mWorkload}, \mref{mPolicy}, \mref{mGA}, \mref{mFindings}, \mref{mDashboard}, \mref{mSuggest}, \mref{mExport}, \mref{mAccess}

\subsection{Syntax}

\subsubsection{Exported Constants}

None.

\subsubsection{Exported Access Programs}

\begin{center}\small
\begin{tabular}{>{\raggedright\arraybackslash}p{3.3cm} >{\raggedright\arraybackslash}p{4.1cm} >{\raggedright\arraybackslash}p{3.4cm} >{\raggedright\arraybackslash}p{4.0cm}}
\hline
\textbf{Name} & \textbf{In} & \textbf{Out} & \textbf{Exceptions} \\
\hline
handle & q: Request & Response & SessionExpired, NotAuthorized \\
\hline
\end{tabular}
\end{center}

\subsection{Semantics}

\subsubsection{State Variables}

None.

\subsubsection{Environment Variables}

None.

\subsubsection{Assumptions}

None.

\subsubsection{Access Routine Semantics}

\todo{Rev 0}

\subsubsection{Local Functions}

\todo{Rev 0}

\newpage


% Harman: MIS of Export, Access Control, Web UI and Request Controller
% (Rev -1: Module, Uses and Syntax).

\section{MIS of Export Module (\mref{mExport})} \label{mis-mExport}

\subsection{Module}

Export

\subsection{Uses}

\mref{mStore}

\subsection{Syntax}

\subsubsection{Exported Constants}

None.

\subsubsection{Exported Access Programs}

\begin{center}\small
\begin{tabular}{>{\raggedright\arraybackslash}p{3.3cm} >{\raggedright\arraybackslash}p{4.1cm} >{\raggedright\arraybackslash}p{3.4cm} >{\raggedright\arraybackslash}p{4.0cm}}
\hline
\textbf{Name} & \textbf{In} & \textbf{Out} & \textbf{Exceptions} \\
\hline
exportData & t: Term, u: User & File & NotAdmin \\
\hline
\end{tabular}
\end{center}

\noindent The file holds only confirmed assessment data, with tentative dates marked.

\subsection{Semantics}

\subsubsection{State Variables}
\todo{Rev 0}

\subsubsection{Environment Variables}
\todo{Rev 0}

\subsubsection{Assumptions}
\todo{Rev 0}

\subsubsection{Access Routine Semantics}
\todo{Rev 0}

\subsubsection{Local Functions}
\todo{Rev 0}

\newpage

\section{MIS of Access Control Module (\mref{mAccess})} \label{mis-mAccess}

\subsection{Module}

AccessControl

\subsection{Uses}

\mref{mStore}

\subsection{Syntax}

\subsubsection{Exported Constants}

\begin{itemize}
\item \textit{SESSION\_TIMEOUT} = 30 min, the time without activity after which a
  session ends (FR\ref{SRS-FR_Access_Timeout} in the SRS~\citep{TeamBVSH2026SRS}).
\end{itemize}

\subsubsection{Exported Access Programs}

\begin{center}\small
\begin{tabular}{>{\raggedright\arraybackslash}p{3.3cm} >{\raggedright\arraybackslash}p{4.1cm} >{\raggedright\arraybackslash}p{3.4cm} >{\raggedright\arraybackslash}p{4.0cm}}
\hline
\textbf{Name} & \textbf{In} & \textbf{Out} & \textbf{Exceptions} \\
\hline
login & cr: Credentials & Session & InvalidCredentials \\
logout & s: Session & -- & -- \\
currentUser & s: Session & User & SessionExpired \\
canEdit & u: User, c: CourseID & $\mathbb{B}$ & -- \\
isAdmin & u: User & $\mathbb{B}$ & -- \\
\hline
\end{tabular}
\end{center}

\subsection{Semantics}

\subsubsection{State Variables}
\todo{Rev 0}

\subsubsection{Environment Variables}
\todo{Rev 0}

\subsubsection{Assumptions}
\todo{Rev 0}

\subsubsection{Access Routine Semantics}
\todo{Rev 0}

\subsubsection{Local Functions}
\todo{Rev 0}

\newpage

\section{MIS of Web UI Module (\mref{mUI})} \label{mis-mUI}

\subsection{Module}

WebUI

\subsection{Uses}

\mref{mControl}

\subsection{Syntax}

\subsubsection{Exported Constants}

None.

\subsubsection{Exported Access Programs}

\begin{center}\small
\begin{tabular}{>{\raggedright\arraybackslash}p{3.3cm} >{\raggedright\arraybackslash}p{4.1cm} >{\raggedright\arraybackslash}p{3.4cm} >{\raggedright\arraybackslash}p{4.0cm}}
\hline
\textbf{Name} & \textbf{In} & \textbf{Out} & \textbf{Exceptions} \\
\hline
showUpload & -- & -- & -- \\
showReview & c: CourseID & -- & -- \\
showCalendar & f: CalendarFilter & -- & -- \\
showFindings & c: CourseID & -- & -- \\
showDashboard & t: Term & -- & -- \\
showExport & -- & -- & -- \\
\hline
\end{tabular}
\end{center}

\noindent Each program displays one page and passes the user's actions on that
page to the Request Controller (\mref{mControl}).

\subsection{Semantics}

\subsubsection{State Variables}
\todo{Rev 0}

\subsubsection{Environment Variables}
\todo{Rev 0}

\subsubsection{Assumptions}
\todo{Rev 0}

\subsubsection{Access Routine Semantics}
\todo{Rev 0}

\subsubsection{Local Functions}
\todo{Rev 0}

\newpage

\section{MIS of Request Controller Module (\mref{mControl})} \label{mis-mControl}

\subsection{Module}

RequestController

\subsection{Uses}

\mref{mIntake}, \mref{mExtract}, \mref{mReview}, \mref{mSchedule}, \mref{mWorkload},
\mref{mPolicy}, \mref{mGA}, \mref{mFindings}, \mref{mDashboard}, \mref{mSuggest},
\mref{mExport}, \mref{mAccess}

\subsection{Syntax}

\subsubsection{Exported Constants}

None.

\subsubsection{Exported Access Programs}

\begin{center}\small
\begin{tabular}{>{\raggedright\arraybackslash}p{3.3cm} >{\raggedright\arraybackslash}p{4.1cm} >{\raggedright\arraybackslash}p{3.4cm} >{\raggedright\arraybackslash}p{4.0cm}}
\hline
\textbf{Name} & \textbf{In} & \textbf{Out} & \textbf{Exceptions} \\
\hline
handle & q: Request & Response & SessionExpired, NotAuthorized \\
\hline
\end{tabular}
\end{center}

\noindent \textit{handle} checks the session and role through Access Control
(\mref{mAccess}) before it calls any other module.

\subsection{Semantics}

\subsubsection{State Variables}
\todo{Rev 0}

\subsubsection{Environment Variables}
\todo{Rev 0}

\subsubsection{Assumptions}
\todo{Rev 0}

\subsubsection{Access Routine Semantics}
\todo{Rev 0}

\subsubsection{Local Functions}
\todo{Rev 0}

\newpage

% Harman: MIS of Cohort and Academic Calendar (Rev -1: Module, Uses and Syntax
% only).

\section{MIS of Cohort Module (\mref{mCohort})} \label{mis-mCohort}

\subsection{Module}

Cohort

\subsection{Uses}

None.

\subsection{Syntax}

\subsubsection{Exported Constants}

None.

\subsubsection{Exported Access Programs}

\begin{center}\small
\begin{tabular}{>{\raggedright\arraybackslash}p{3.4cm} >{\raggedright\arraybackslash}p{3.3cm} >{\raggedright\arraybackslash}p{4.1cm} >{\raggedright\arraybackslash}p{3.8cm}}
\hline
\textbf{Name} & \textbf{In} & \textbf{Out} & \textbf{Exceptions} \\
\hline
isDepartmentCourse & k: CourseCode & $\mathbb{B}$ & -- \\
cohortOf & k: CourseCode & Cohort & UnknownCode \\
coursesIn & h: Cohort & set of CourseCode & UnknownCohort \\
\hline
\end{tabular}
\end{center}

\noindent A cohort is a (program, year level) pair, for example (COMPSCI, 2).
Only department courses on the master course list belong to a cohort.

\subsection{Semantics}

\subsubsection{State Variables}
\todo{Rev 0}

\subsubsection{Environment Variables}
\todo{Rev 0}

\subsubsection{Assumptions}
\todo{Rev 0}

\subsubsection{Access Routine Semantics}
\todo{Rev 0}

\subsubsection{Local Functions}
\todo{Rev 0}

\newpage

\section{MIS of Academic Calendar Module (\mref{mCalendar})} \label{mis-mCalendar}

\subsection{Module}

AcademicCalendar

\subsection{Uses}

None.

\subsection{Syntax}

\subsubsection{Exported Constants}

None.

\subsubsection{Exported Access Programs}

\begin{center}\small
\begin{tabular}{>{\raggedright\arraybackslash}p{3.4cm} >{\raggedright\arraybackslash}p{3.3cm} >{\raggedright\arraybackslash}p{4.1cm} >{\raggedright\arraybackslash}p{3.8cm}}
\hline
\textbf{Name} & \textbf{In} & \textbf{Out} & \textbf{Exceptions} \\
\hline
termDates & t: Term & (start: Date, end: Date) & UnknownTerm \\
classDays & t: Term & seq of Date & UnknownTerm \\
weekOf & d: Date, t: Term & $\mathbb{N}$ & UnknownTerm, DateOutsideTerm \\
isRestricted & d: Date & $\mathbb{B}$ & DateOutsideTerm \\
restrictedPeriods & t: Term & seq of (start: Date, end: Date, kind: String) & UnknownTerm \\
\hline
\end{tabular}
\end{center}

\noindent \textit{isRestricted} is true when the date falls in a holiday, a
recess or the examination period. The ban period is not included: Policy
Compliance (\mref{mPolicy}) works it out from \textit{classDays}.

\subsection{Semantics}

\subsubsection{State Variables}
\todo{Rev 0}

\subsubsection{Environment Variables}
\todo{Rev 0}

\subsubsection{Assumptions}
\todo{Rev 0}

\subsubsection{Access Routine Semantics}
\todo{Rev 0}

\subsubsection{Local Functions}
\todo{Rev 0}

\newpage

% Harman: MIS of the software decision modules (Rev -1: Module, Uses and
% Syntax only).

\section{MIS of Document Reader Module (\mref{mReader})} \label{mis-mReader}

\subsection{Module}

DocumentReader

\subsection{Uses}

Hardware-Hiding (\mref{mHH})

\subsection{Syntax}

\subsubsection{Exported Constants}

None.

\subsubsection{Exported Access Programs}

\begin{center}\small
\begin{tabular}{>{\raggedright\arraybackslash}p{2.2cm} >{\raggedright\arraybackslash}p{4.2cm} >{\raggedright\arraybackslash}p{4.3cm} >{\raggedright\arraybackslash}p{3.9cm}}
\hline
\textbf{Name} & \textbf{In} & \textbf{Out} & \textbf{Exceptions} \\
\hline
readText & f: File & seq of (page: $\mathbb{N}$, text: String) & UnreadableDocument, UnsupportedFormat \\
highlight & f: File, s: seq of (page: $\mathbb{N}$, text: String) & (File, seq of (page: $\mathbb{N}$, text: String)) & UnreadableDocument, UnsupportedFormat \\
\hline
\end{tabular}
\end{center}

\noindent The output of \textit{highlight} is the highlighted copy of the
outline and the passages in \textit{s} that could not be found on their page.

\subsection{Semantics}

\subsubsection{State Variables}
\todo{Rev 0}

\subsubsection{Environment Variables}
\todo{Rev 0}

\subsubsection{Assumptions}
\todo{Rev 0}

\subsubsection{Access Routine Semantics}
\todo{Rev 0}

\subsubsection{Local Functions}
\todo{Rev 0}

\newpage

\section{MIS of LLM Client Module (\mref{mLLM})} \label{mis-mLLM}

\subsection{Module}

LLMClient

\subsection{Uses}

Hardware-Hiding (\mref{mHH})

\subsection{Syntax}

\subsubsection{Exported Constants}

None.

\subsubsection{Exported Access Programs}

\begin{center}\small
\begin{tabular}{>{\raggedright\arraybackslash}p{2.2cm} >{\raggedright\arraybackslash}p{4.2cm} >{\raggedright\arraybackslash}p{4.3cm} >{\raggedright\arraybackslash}p{3.9cm}}
\hline
\textbf{Name} & \textbf{In} & \textbf{Out} & \textbf{Exceptions} \\
\hline
complete & i: String, d: String & String & ModelUnavailable, Timeout \\
isAvailable & -- & $\mathbb{B}$ & -- \\
\hline
\end{tabular}
\end{center}

\noindent In \textit{complete}, \textit{i} holds the instructions and
\textit{d} holds the outline text. They are kept apart so that outline text is
always passed to the model as data, never as instructions.

\subsection{Semantics}

\subsubsection{State Variables}
\todo{Rev 0}

\subsubsection{Environment Variables}
\todo{Rev 0}

\subsubsection{Assumptions}
\todo{Rev 0}

\subsubsection{Access Routine Semantics}
\todo{Rev 0}

\subsubsection{Local Functions}
\todo{Rev 0}

\newpage

\section{MIS of Data Store Module (\mref{mStore})} \label{mis-mStore}

\subsection{Module}

DataStore

\subsection{Uses}

Hardware-Hiding (\mref{mHH})

\subsection{Syntax}

\subsubsection{Exported Constants}

\begin{itemize}
\item \textit{RETENTION\_YEARS} = 10, the retention period for outlines,
  extracted data and findings, from the data requirements of the
  SRS~\citep{TeamBVSH2026SRS}.
\end{itemize}

\subsubsection{Exported Access Programs}

\begin{center}\small
\begin{tabular}{>{\raggedright\arraybackslash}p{2.2cm} >{\raggedright\arraybackslash}p{4.2cm} >{\raggedright\arraybackslash}p{4.3cm} >{\raggedright\arraybackslash}p{3.9cm}}
\hline
\textbf{Name} & \textbf{In} & \textbf{Out} & \textbf{Exceptions} \\
\hline
save & e: Entity & ID & InvalidEntity \\
load & k: EntityKind, id: ID & Entity & UnknownKind, NotFound \\
query & k: EntityKind, p: Filter & seq of Entity & UnknownKind \\
saveFile & f: File & ID & -- \\
loadFile & id: ID & File & NotFound \\
\hline
\end{tabular}
\end{center}

\noindent \textit{save} adds a new entity or replaces a stored one with the
same ID. There is no delete program: records are kept for
\textit{RETENTION\_YEARS}.

\subsection{Semantics}

\subsubsection{State Variables}
\todo{Rev 0}

\subsubsection{Environment Variables}
\todo{Rev 0}

\subsubsection{Assumptions}
\todo{Rev 0}

\subsubsection{Access Routine Semantics}
\todo{Rev 0}

\subsubsection{Local Functions}
\todo{Rev 0}

\newpage

% Same IEEE style file as the SRS (IEEEtran.bst is not installed in CI).
\bibliographystyle {../../SRS-IEEE/IEEEtran}
\bibliography {../../../refs/References}

\newpage

\newpage{}

\section{Appendix --- Generative AI Use Disclosure}

\wss{ChatGPT (version 5) was used to brainstorm ideas for the template for this disclose section.}

\subsection{AI Tools Used}

\wss{Give the name of the AI tool(s) used and the version.}

\subsection{How and Where Used}

\wss{Explain what the tool(s) were used for.  Examples include: brainstorming,
explaining a technical concept, reviewing the document, generating or modifying
text, generating or modifying diagrams, identifying errors or omissions or
inconsistencies.}

\wss{For each of the uses, identify which parts of the document were affected.}

\wss{You do not need to report every interaction, but you should disclose the 
uses that materially contributed to your work}

\subsection{Verification of AI-Generated Content}

\wss{\begin{itemize}
    \item Describe how the team checked AI-generated or AI-assisted material
    before including it in the document.
    \item In particular, verify factual claims, technical assertions,
    requirements, calculations, references, and other information for which an
    AI system may produce incorrect or fabricated results.
    \item \textbf{The team is responsible for the accuracy, completeness,
consistency, and appropriateness of the submitted document, regardless of
whether material was generated by a human or an AI system.}
\end{itemize}}

\newpage{}

\section{Appendix --- Reflection}

\wss{Not required for CAS 741 projects}

The information in this section will be used to evaluate the team members on the
graduate attribute of Problem Analysis and Design.

\input{../../Reflection.text}

\begin{enumerate}
  \item What went well while writing this deliverable? 
  \item What pain points did you experience during this deliverable, and how
    did you resolve them?
  \item Which of your design decisions stemmed from speaking to your client(s)
  or a proxy (e.g. your peers, stakeholders, potential users)? For those that
  were not, why, and where did they come from?
  \item While creating the design doc, what parts of your other documents (e.g.
  requirements, hazard analysis, etc), it any, needed to be changed, and why?
  \item What are the limitations of your solution?  Put another way, given
  unlimited resources, what could you do to make the project better? (LO\_ProbSolutions)
  \item Give a brief overview of other design solutions you considered.  What
  are the benefits and tradeoffs of those other designs compared with the chosen
  design?  From all the potential options, why did you select the documented design?
  (LO\_Explores)
  \item (After you have implemented another team's module, which means this
  isn't filled in until after the original deadline). What did you learn by
  implementing another team's module? Were all the details you needed in the
  documentation, or did you need to make assumptions, or ask the other team
  questions?  If your team also had another team implement one of your modules,
  what was this experience like?  Are there things in your documentation you
  could have changed to make the process go more smoothly for when an
  ``outsider'' completes some of the implementation?

  Your reflection should be about concepts learned in your SE education, such as modular decomposition, cohesion, coupling, information hiding, etc.  You should critique the quality and utility of the specification you received.

  You should create a PR for your implemented module.  A link to the PR should be part of your reflection.

\end{enumerate}

\end{document}